% Generated draft: requires mathematical and source review.
\documentclass[11pt]{article}
\usepackage[margin=25mm]{geometry}
\usepackage{fontspec}
\setmainfont{Latin Modern Roman}
\usepackage{amsmath,amssymb,mathtools,unicode-math}
\setmathfont{Latin Modern Math}
\usepackage{enumitem,needspace,xurl,hyperref}
\setlength{\parindent}{0pt}
\setlength{\parskip}{5pt}
\setlength{\emergencystretch}{3em}
\newcommand{\N}{\mathbb{N}}
\newcommand{\Z}{\mathbb{Z}}
\newcommand{\Q}{\mathbb{Q}}
\newcommand{\R}{\mathbb{R}}
\newcommand{\C}{\mathbb{C}}
\newcommand{\F}{\mathbb{F}}
\newcommand{\A}{\mathbb{A}}
\newcommand{\PP}{\mathbb{P}}
\newcommand{\E}{\mathbb{E}}
\newcommand{\eps}{\varepsilon}
\newcommand{\dd}{\,\mathrm{d}}
\begin{document}
\section*{UnsolvedMath ID: 30005515}
\textbf{Source problem code:} catalogue entry
\subsection{Definitions and mathematical statement}
Let \(p\) be a prime and let \(k\) be a field of characteristic \(p\). For a finite group \(G\), its group algebra is
\[
 kG=\left\{\sum_{g\in G}a_g g:a_g\in k\right\},
\]
with coefficientwise addition and multiplication induced by the group operation and extended \(k\)-bilinearly. All group-algebra isomorphisms below mean isomorphisms of unital \(k\)-algebras; group isomorphisms are indicated separately. The prime subfield of \(k\) is identified with \(\F_p\).

For a class \(\mathcal C\) of finite \(p\)-groups, define the recovery assertion
\[
\begin{split}
\mathcal M_p(\mathcal C):\quad
&\text{for every field }k\text{ of characteristic }p,\\
&\text{every }G\in\mathcal C\text{ and every finite }p\text{-group }H,\\
&kG\cong_{k\text{-alg}}kH\ \Longrightarrow\ G\cong_{\mathrm{grp}}H.
\end{split}
\]
This makes explicit a usual finite \(p\)-group formulation of the Modular Isomorphism Problem. The supplied excerpt does not give its preceding definition. Consequently, the field quantifier, the restriction to finite \(p\)-groups, and whether \(H\) must also belong to the specified subclass require confirmation from the complete source context. The first two questions below are stated under this explicit, provisional interpretation.

Write \([G,G]\) for the subgroup generated by the commutators \([x,y]=x^{-1}y^{-1}xy\), and \(Z(G)\) for the center. A group has nilpotency class exactly \(2\) when
\[
1\ne [G,G]\subseteq Z(G).
\]
Let \(\mathcal P_p\) be the class of finite \(p\)-groups and let \(\mathcal C_{p,2}\subseteq\mathcal P_p\) consist of those of nilpotency class exactly \(2\).

\begin{enumerate}
\item Determine whether \(\mathcal M_p(\mathcal P_p)\) holds for every odd prime \(p\). This expresses the odd-order question in the finite \(p\)-group setting. The literal phrase ``groups of odd order'' could instead include finite groups whose orders have several prime divisors; the supplied record does not specify that broader formulation.
\item Determine whether \(\mathcal M_p(\mathcal C_{p,2})\) holds for every prime \(p\). The source says class \(2\); whether it intends exactly \(2\) or at most \(2\) is not explicit.
\item Determine whether there exist a prime \(p\), finite \(p\)-groups \(G,H\), and a field \(k\) of characteristic \(p\) such that
\[
\F_pG\not\cong_{\F_p\text{-alg}}\F_pH,
\qquad
kG\cong_{k\text{-alg}}kH.
\]
There is no supplied restriction that \(k\) be finite or algebraic over \(\F_p\).
\item Determine which group invariants are determined by the group algebra. More precisely, for a fixed field \(k\) of characteristic \(p\), an isomorphism-invariant function \(I\) on finite \(p\)-groups is determined by \(kG\) if
\[
\forall G,H\in\mathcal P_p,\qquad
kG\cong_{k\text{-alg}}kH\ \Longrightarrow\ I(G)=I(H).
\]
A group property is covered by taking \(I\) to be its indicator function. The request is to identify such invariants and properties, including their dependence on \(k\). The source supplies neither a prescribed family of invariants nor a precise completeness criterion for this classification request; its intended group domain also needs confirmation.
\end{enumerate}

The original excerpt contains the self-referential sentence that the isomorphism type of \(kG\) cannot be recovered from \(kG\). Taken literally, that sentence does not express a coherent obstruction to recovery. The supplied corrected statement instead asks which properties of \(G\) are determined by \(kG\); item~4 follows that corrected statement while retaining the textual discrepancy for review.
\subsection{Short English statement}
The record asks whether modular group algebras recover odd-order groups and groups of nilpotency class \(2\), whether extending the coefficient field can make two previously nonisomorphic group algebras isomorphic, and which group invariants the algebra determines. The finite \(p\)-group scope and quantifiers of the first two questions are provisional because their defining context is absent. This is a statement-only draft requiring human mathematical review.
\subsection{Sources}
\begin{itemize}
\item Local registry: \texttt{data/unsolvedmath/problems.json}, with \texttt{catalog/problems.json} as the additional-record registry; permanent UnsolvedMath ID \(30005515\), title ``Four Modular Isomorphism Problems,'' exact source code \texttt{catalogue entry}. These registry identifiers are supplied project provenance; their files were not inspected for this draft.
\item Supplied source report: \emph{Representations of Finite Groups} (2023), Oberwolfach Reports, pp.~1031--1084, DOI \texttt{10.4171/owr/2023/19}. \url{https://doi.org/10.4171/owr/2023/19}.
\item Supplied attribution: ``Modular Isomorphism Problem -- progress, solution and open challenges,'' Leo Margolis, joint work with D.~Garc\'{i}a-Lucas, \'{A}.~del R\'{i}o, T.~Sakurai, and M.~Stanojkovski. The record gives proposed year \(2023\) and no separately named proposer.
\item Supplied literature source, labeled ``modular isomorphism survey'': \url{https://doi.org/10.1365/s13291-022-00249-5}.
\item Supplied literature source, labeled ``class-two cyclic-center result'': \url{https://doi.org/10.1515/forum-2023-0237}.
\item Literature assessment dated \texttt{2026-08-22}: the supplied curation labels the problem open, describes the general odd-order and class-two questions as remaining open, and describes subclasses involving cyclic center and center index \(p^3\) as solved. These are supplied dated curation claims, not independently verified mathematical or literature conclusions. The hypotheses of those subclass results are not fully stated in the record and require independent checking.
\item The registry labels the statement \texttt{corrected\_verified} and reports that it was reconstructed and checked against the official Oberwolfach report. This is a supplied verification claim; no source was visited or independently checked in preparing this draft.
\end{itemize}
\end{document}
